\documentclass[aspectratio=169, 10pt]{beamer}

\usetheme{Madrid}
\usefonttheme{professionalfonts}

\usepackage[utf8]{inputenc}
\usepackage[indonesian]{babel}
\usepackage{amsmath,amsfonts,amssymb}
\usepackage{graphicx}
\usepackage{booktabs}
\usepackage{listings}
\usepackage{xcolor}

\definecolor{mainblue}{RGB}{24, 43, 73}
\definecolor{accentblue}{RGB}{41, 128, 185}
\definecolor{accentgold}{RGB}{230, 126, 34}
\definecolor{darkgray}{RGB}{44, 62, 80}
\definecolor{lightbg}{RGB}{248, 250, 252}
\definecolor{bordergray}{RGB}{200, 210, 220}
\definecolor{codebg}{RGB}{240, 243, 246}
\definecolor{codegreen}{RGB}{39, 174, 96}

\setbeamercolor{structure}{fg=mainblue}
\setbeamercolor{palette primary}{bg=mainblue, fg=white}
\setbeamercolor{palette secondary}{bg=accentblue, fg=white}
\setbeamercolor{palette tertiary}{bg=mainblue!80, fg=white}
\setbeamercolor{title}{fg=white, bg=mainblue}
\setbeamercolor{frametitle}{fg=white, bg=mainblue}
\setbeamercolor{block title}{bg=mainblue, fg=white}
\setbeamercolor{block body}{bg=lightbg, fg=darkgray}
\setbeamercolor{block title alerted}{bg=accentgold, fg=white}
\setbeamercolor{block body alerted}{bg=accentgold!10, fg=darkgray}

\lstset{
    language=Python,
    backgroundcolor=\color{codebg},
    basicstyle=\ttfamily\scriptsize\color{darkgray},
    keywordstyle=\bfseries\color{accentblue},
    stringstyle=\color{accentgold},
    commentstyle=\itshape\color{codegreen},
    showstringspaces=false,
    breaklines=true,
    frame=single,
    rulecolor=\color{bordergray},
    numbers=none
}

\title[\textbf{Deret Waktu (Minggu 6)}]{\textbf{Pemodelan Eksogen \& Multivariat: SARIMAX \& VAR}}
\subtitle{Korelasi Silang (CCF), Pengaruh Cuaca pada Beban Listrik, dan Sistem Multi-Kanal}
\author[\textbf{Ir. Novalio Daratha, Ph.D.}]{\textbf{Ir. Novalio Daratha, S.T., M.Sc., Ph.D.}}
\institute[FT UNIB]{Program Studi Teknik Elektro -- Fakultas Teknik\\Universitas Bengkulu}
\date{Minggu 6 / Semester Ganjil 2026}

\begin{document}

\begin{frame}[plain]
    \titlepage
\end{frame}

\begin{frame}{Agenda Perkuliahan Minggu ke-6}
    \begin{columns}[T]
        \begin{column}{0.48\textwidth}
            \begin{block}{Capaian Pembelajaran (Sub-CPMK 6)}
                Mahasiswa mampu mengintegrasikan variabel luar (eksogen) seperti cuaca/temperatur ke dalam model SARIMAX, menganalisis korelasi silang (CCF), serta memformulasikan sistem multivariat Vector AutoRegression (VAR).
            \end{block}
        \end{column}
        \begin{column}{0.48\textwidth}
            \begin{enumerate}
                \item \textbf{Urgensi Variabel Eksogen di Sistem Tenaga}
                \item \textbf{Fungsi Korelasi Silang (CCF)}
                \item \textbf{Formulasi Model SARIMAX}
                \item \textbf{Model Transfer Function dan Distributed Lag}
                \item \textbf{Pengantar Vector AutoRegression (VAR)}
                \item \textbf{Implementasi Python: Beban vs Suhu Udara}
            \end{enumerate}
        \end{column}
    \end{columns}
\end{frame}

\begin{frame}{Mengapa Membutuhkan Variabel Eksogen ($X_t$)?}
    \begin{columns}[T]
        \begin{column}{0.5\textwidth}
            \begin{block}{Faktor Fisik Pembentuk Beban Listrik}
                \begin{itemize}
                    \item \textbf{Temperatur Udara}: Saat cuaca panas, penggunaan AC pendingin ruangan melonjak drastis (\textit{Cooling Degree Days}).
                    \item \textbf{Radiasi Matahari}: Mempengaruhi keluaran daya fotovoltaik (PV) di atap gedung.
                    \item \textbf{Kalender \& Hari Libur}: Perbedaan pola hari kerja, akhir pekan, dan hari libur nasional.
                \end{itemize}
            \end{block}
        \end{column}
        \begin{column}{0.5\textwidth}
            \begin{alertblock}{Keterbatasan SARIMA Murni}
                SARIMA univariat hanya berasumsi bahwa masa depan dibentuk oleh masa lalu sinyal itu sendiri ($y_t = f(y_{t-k})$). Jika ada gelombang panas (*heatwave*) mendadak, SARIMA murni akan terlambat merespons.
            \end{alertblock}
        \end{column}
    \end{columns}
\end{frame}

\begin{frame}{Cross-Correlation Function (CCF)}
    Sebelum memasukkan variabel eksogen, kita menganalisis keterlambatan efek (*lagged relationship*) menggunakan CCF:
    \begin{block}{Definisi Matematis CCF}
        \begin{equation}
            r_{xy}(k) = \frac{\sum_{t=1}^{N-k} (x_t - \bar{x})(y_{t+k} - \bar{y})}{\sqrt{\sum_{t=1}^N (x_t - \bar{x})^2 \sum_{t=1}^N (y_t - \bar{y})^2}}
        \end{equation}
    \end{block}
    \begin{columns}[T]
        \begin{column}{0.48\textwidth}
            \textbf{Interpretasi Lag Positif ($k > 0$):}
            Perubahan suhu pada waktu $t$ baru mempengaruhi lonjakan konsumsi beban AC pada waktu $t+k$ (efek inersia termal bangunan, misal $k=2$ jam).
        \end{column}
        \begin{column}{0.48\textwidth}
            \textbf{Pra-syarat Prewhitening:}
            Kedua sinyal harus di-\textit{prewhiten} agar korelasi tidak palsu (*spurious correlation*) akibat tren bersama.
        \end{column}
    \end{columns}
\end{frame}

\begin{frame}{Formulasi Matematis SARIMAX}
    Menggabungkan regresi linier dengan derau terstruktur SARIMA:
    \begin{block}{Persamaan SARIMAX}
        \begin{equation}
            y_t = \sum_{j=1}^m \beta_j X_{j,t} + \eta_t
        \end{equation}
        di mana galat $\eta_t$ mengikuti proses SARIMA musiman:
        \begin{equation}
            \Phi_P(B^s)\phi_p(B)(1 - B)^d(1 - B^s)^D \eta_t = \Theta_Q(B^s)\theta_q(B)\epsilon_t
        \end{equation}
    \end{block}
    \begin{alertblock}{Keunggulan}
        Koefisien $\beta_j$ merefleksikan elastisitas fisik: misal kenaikan $1^\circ\text{C}$ suhu meningkatkan beban sebesar $15.4$ MW.
    \end{alertblock}
\end{frame}

\begin{frame}{Pengantar Vector AutoRegression (VAR)}
    Jika terdapat beberapa variabel kelistrikan yang saling mempengaruhi secara timbal balik (misal Tegangan Bus $V_1, V_2, V_3$ pada gardu transmisi):
    \begin{block}{Persamaan Model VAR($p$)}
        \begin{equation}
            \mathbf{Y}_t = \mathbf{c} + \mathbf{A}_1 \mathbf{Y}_{t-1} + \mathbf{A}_2 \mathbf{Y}_{t-2} + \dots + \mathbf{A}_p \mathbf{Y}_{t-p} + \boldsymbol{\epsilon}_t
        \end{equation}
        di mana $\mathbf{Y}_t = [y_{1,t}, y_{2,t}, \dots, y_{k,t}]^T$ adalah vektor observasi $k$-dimensi, dan $\mathbf{A}_i$ adalah matriks koefisien $k \times k$.
    \end{block}
    \vspace{0.1cm}
    \begin{center}
        \textbf{Uji Kausalitas Granger}: Menguji apakah riwayat sinyal $X$ secara statistik membantu meramalkan sinyal $Y$.
    \end{center}
\end{frame}

\begin{frame}[fragile]{Implementasi Python: SARIMAX dengan Data Suhu}
\begin{lstlisting}
from statsmodels.tsa.statespace.sarimax import SARIMAX
import pandas as pd

# Load data beban & temperatur
df = pd.read_csv('load_weather.csv', parse_dates=['time'], index_col='time')
train, test = df[:-24], df[-24:]

# Definisikan variabel eksogen (X: suhu, radiasi)
exog_train = train[['temperature', 'solar_irradiance']]
exog_test = test[['temperature', 'solar_irradiance']]

# Fitting Model SARIMAX(1,1,1)(1,1,0)_24
sarimax = SARIMAX(train['load_mw'], exog=exog_train,
                  order=(1, 1, 1), seasonal_order=(1, 1, 0, 24)).fit(disp=False)

print(sarimax.summary())

# Forecasting 24 Jam ke Depan dengan Ramalan Cuaca
pred_24 = sarimax.forecast(steps=24, exog=exog_test)
\end{lstlisting}
\end{frame}

\begin{frame}[fragile]{Implementasi Python: Vector AutoRegression (VAR)}
\begin{lstlisting}
from statsmodels.tsa.api import VAR
import pandas as pd
import matplotlib.pyplot as plt

# Load dataset sistem 3-variabel (Tegangan, Arus, THD)
df_multi = pd.read_csv('power_quality.csv', index_col='time')

# Inisialisasi dan pencarian ordo optimal VAR (minimasi AIC)
model_var = VAR(df_multi)
results = model_var.fit(maxlags=15, ic='aic')
print(f"Ordo VAR terpilih: {results.k_ar}")
print(results.summary())

# Analisis Impulse Response Function (IRF)
# Melihat respon variabel lain jika ada lonjakan sesaat di 'Arus'
irf = results.irf(10) # 10 langkah ke depan
irf.plot(impulse='Arus', response='THD')
plt.title("Respon THD terhadap Impuls Arus")
plt.show()

# Peramalan sistem multivariat
forecast_var = results.forecast(df_multi.values[-results.k_ar:], steps=5)
\end{lstlisting}
\end{frame}


\begin{frame}{Rangkuman \& Persiapan Minggu ke-7}
    \begin{block}{Poin Kunci Minggu 6}
        \begin{itemize}
            \item SARIMAX memperluas kemampuan peramalan dengan menyerap informasi kausal fisik eksternal.
            \item Analisis CCF mendeteksi jeda waktu inersia termal sistem kelistrikan.
            \item VAR memodelkan dinamika interaksi multi-bus secara simultan tanpa asumsi searah.
        \end{itemize}
    \end{block}
    \vspace{0.2cm}
    \begin{alertblock}{Materi Minggu Depan (Minggu ke-7)}
        \textbf{Estimasi State \& Kalman Filter}: Representasi ruang keadaan (\textit{State-Space}), algoritma Discrete Kalman Filter, dan peredaman derau pengukuran sensor telemetri.
    \end{alertblock}
\end{frame}

\begin{frame}[plain]
    \begin{center}
        \vspace{1.5cm}
        {\Huge \color{mainblue} \textbf{Terima Kasih}}\\[0.4cm]
        {\large Sesi Tanya Jawab \& Diskusi Model Eksogen}\\[1.2cm]
        \textbf{Laboratorium Sistem Tenaga \& Komputasi Cerdas}\\
        Jurusan Teknik Elektro -- Universitas Bengkulu
    \end{center}
\end{frame}

\end{document}
